\section{Worked solutions on the ten soils (second scenario)}
The same questions are repeated for different working conditions, to show how the answer moves with them.

\exbox{Ten-soil example 13b: Coarse loam (S01)}{Coarse loam ($K_s=13.0$ mm\,h$^{-1}$): a 4.0-mm nozzle at 300 kPa works on a $12\times18$ m spacing. Find the discharge, the application rate and whether runoff will occur.}\sol $h=350/9.81=30.6$ m; $q=0.95\pi(30.6/2000)^2\sqrt{2\times9.81\times4.0}\times6\times10^4=17.5$ L\,min$^{-1}$ \ans $r=q\times60/(S_lS_m)=17.5\times60/216=4.9$ mm\,h$^{-1}$ \ans The rate is below $K_s$, so no runoff. \ans

\exbox{Ten-soil example 13b: Coarse sand (S02)}{Coarse sand ($K_s=120.0$ mm\,h$^{-1}$): a 4.0-mm nozzle at 300 kPa works on a $12\times18$ m spacing. Find the discharge, the application rate and whether runoff will occur.}\sol $h=350/9.81=30.6$ m; $q=0.95\pi(30.6/2000)^2\sqrt{2\times9.81\times4.0}\times6\times10^4=17.5$ L\,min$^{-1}$ \ans $r=q\times60/(S_lS_m)=17.5\times60/216=4.9$ mm\,h$^{-1}$ \ans The rate is below $K_s$, so no runoff. \ans

\exbox{Ten-soil example 13b: Sand (S03)}{Sand ($K_s=210.0$ mm\,h$^{-1}$): a 4.0-mm nozzle at 300 kPa works on a $12\times18$ m spacing. Find the discharge, the application rate and whether runoff will occur.}\sol $h=350/9.81=30.6$ m; $q=0.95\pi(30.6/2000)^2\sqrt{2\times9.81\times4.0}\times6\times10^4=17.5$ L\,min$^{-1}$ \ans $r=q\times60/(S_lS_m)=17.5\times60/216=4.9$ mm\,h$^{-1}$ \ans The rate is below $K_s$, so no runoff. \ans

\exbox{Ten-soil example 13b: Loamy sand (S04)}{Loamy sand ($K_s=61.0$ mm\,h$^{-1}$): a 4.0-mm nozzle at 300 kPa works on a $12\times18$ m spacing. Find the discharge, the application rate and whether runoff will occur.}\sol $h=350/9.81=30.6$ m; $q=0.95\pi(30.6/2000)^2\sqrt{2\times9.81\times4.0}\times6\times10^4=17.5$ L\,min$^{-1}$ \ans $r=q\times60/(S_lS_m)=17.5\times60/216=4.9$ mm\,h$^{-1}$ \ans The rate is below $K_s$, so no runoff. \ans

\exbox{Ten-soil example 13b: Sandy loam (S05)}{Sandy loam ($K_s=26.0$ mm\,h$^{-1}$): a 4.0-mm nozzle at 300 kPa works on a $12\times18$ m spacing. Find the discharge, the application rate and whether runoff will occur.}\sol $h=350/9.81=30.6$ m; $q=0.95\pi(30.6/2000)^2\sqrt{2\times9.81\times4.0}\times6\times10^4=17.5$ L\,min$^{-1}$ \ans $r=q\times60/(S_lS_m)=17.5\times60/216=4.9$ mm\,h$^{-1}$ \ans The rate is below $K_s$, so no runoff. \ans

\exbox{Ten-soil example 13b: Loam (S06)}{Loam ($K_s=13.0$ mm\,h$^{-1}$): a 4.0-mm nozzle at 300 kPa works on a $12\times18$ m spacing. Find the discharge, the application rate and whether runoff will occur.}\sol $h=350/9.81=30.6$ m; $q=0.95\pi(30.6/2000)^2\sqrt{2\times9.81\times4.0}\times6\times10^4=17.5$ L\,min$^{-1}$ \ans $r=q\times60/(S_lS_m)=17.5\times60/216=4.9$ mm\,h$^{-1}$ \ans The rate is below $K_s$, so no runoff. \ans

\exbox{Ten-soil example 13b: Silt loam (S07)}{Silt loam ($K_s=6.8$ mm\,h$^{-1}$): a 4.0-mm nozzle at 300 kPa works on a $12\times18$ m spacing. Find the discharge, the application rate and whether runoff will occur.}\sol $h=350/9.81=30.6$ m; $q=0.95\pi(30.6/2000)^2\sqrt{2\times9.81\times4.0}\times6\times10^4=17.5$ L\,min$^{-1}$ \ans $r=q\times60/(S_lS_m)=17.5\times60/216=4.9$ mm\,h$^{-1}$ \ans The rate is below $K_s$, so no runoff. \ans

\exbox{Ten-soil example 13b: Sandy clay loam (S08)}{Sandy clay loam ($K_s=4.3$ mm\,h$^{-1}$): a 4.0-mm nozzle at 300 kPa works on a $12\times18$ m spacing. Find the discharge, the application rate and whether runoff will occur.}\sol $h=350/9.81=30.6$ m; $q=0.95\pi(30.6/2000)^2\sqrt{2\times9.81\times4.0}\times6\times10^4=17.5$ L\,min$^{-1}$ \ans $r=q\times60/(S_lS_m)=17.5\times60/216=4.9$ mm\,h$^{-1}$ \ans The rate exceeds $K_s$: ponding after $101.0$ h; sets longer than that run off, so use a smaller nozzle or a wider spacing. \ans

\exbox{Ten-soil example 13b: Clay loam (S09)}{Clay loam ($K_s=2.3$ mm\,h$^{-1}$): a 4.0-mm nozzle at 300 kPa works on a $12\times18$ m spacing. Find the discharge, the application rate and whether runoff will occur.}\sol $h=350/9.81=30.6$ m; $q=0.95\pi(30.6/2000)^2\sqrt{2\times9.81\times4.0}\times6\times10^4=17.5$ L\,min$^{-1}$ \ans $r=q\times60/(S_lS_m)=17.5\times60/216=4.9$ mm\,h$^{-1}$ \ans The rate exceeds $K_s$: ponding after $11.9$ h; sets longer than that run off, so use a smaller nozzle or a wider spacing. \ans

\exbox{Ten-soil example 13b: Silty clay loam (S10)}{Silty clay loam ($K_s=1.5$ mm\,h$^{-1}$): a 4.0-mm nozzle at 300 kPa works on a $12\times18$ m spacing. Find the discharge, the application rate and whether runoff will occur.}\sol $h=350/9.81=30.6$ m; $q=0.95\pi(30.6/2000)^2\sqrt{2\times9.81\times4.0}\times6\times10^4=17.5$ L\,min$^{-1}$ \ans $r=q\times60/(S_lS_m)=17.5\times60/216=4.9$ mm\,h$^{-1}$ \ans The rate exceeds $K_s$: ponding after $7.5$ h; sets longer than that run off, so use a smaller nozzle or a wider spacing. \ans

